\section{The fourth root on each component and the different at a finite prime}
\label{sec:key-different}

The weight-one twist of Section~\ref{sec:jet-connection} comes from a
weight-four theta rule. Here we explain its descent from the complex
chart and compute the precise determinant loss in a vertical
valuation. Write \(k\) for the coefficient number field and
\(m_\tau\) for the multiplicity of the embedding \(\tau:E\hookrightarrow k\)
in the Lie algebra. Repetitions of a single embedding occur through
the integer \(m_\tau\). Embeddings of multiplicity zero may be retained.

\subsection{Descending the fourth root}

On each smooth framed chart with coordinate ring \(A\), the
functorial theta scalar in a primitive frame is a unit \(u\in A^\times\).
Its transformation law under a triangular change of frame is the
fourth power of the frame character. The chartwise values therefore
give the weight-four rule \(K_4\), and the reciprocal values glue
in the opposite weight. Their product has value one at each generic
component, hence is the identity. The complex uniformization gives
one more piece of data: after embedding \(k\) in \(\mathbb C\), the
complex-valued function \(C u\), \(C\in\mathbb C^\times\), has a
continuous fourth root on the complex points of the chart. The root
is built from the determinant of a matrix of holomorphic periods,
divided by the square of the scalar relating the primitive theta
frames, with the common normalization from the Hermitian datum.
The adapted-period charts and the transformation law of the normalized
frames make this value independent of the representative and
continuous. These are hypotheses on the actual normalized theta
uniformization. A fourth root at individual complex points would
give no descent.

Here is the algebraic descent. Put
\(B=\mathbb C\otimes_k A\), a smooth integral affine algebra, and
\(u'=Cu\) in \(B^\times\). The integral-domain assertion here is
the geometric integrality of the chart of the chosen component.
The finite cover \(t^4=u'\) is étale, since \(4t^3\) is invertible.
On any connected component of its complex points the ratio of \(t\)
to the supplied continuous fourth root belongs to \(\mu_4\), so it
is constant. An integral algebraic component has connected complex
points. Smoothness makes the algebraic components disjoint, and
each maps onto the connected base, as the image of a finite étale
component is both open and closed. For one such component the
constant-ratio assertion gives at most one point above each complex
point of the base. Its rank is consequently one, and the component
maps isomorphically to the base. This gives \(b'\in B^\times\)
with \(b'^4=u'\).

Choose a fourth root of \(C\) in \(\mathbb C\) and absorb it into
\(b'\), so that \(b'^4=u\). Each automorphism of \(\mathbb C/k\)
takes \(b'\) to a fourth root of the same unit. The ratio is a
constant fourth root of unity: the polynomial \(t^4-1\) splits into
linear factors over \(\mathbb C\), and \(B\) is a domain. Expand
\(b'=\sum_j c_j a_j\) in finitely many members of a \(k\)-basis of
\(A\), with \(c_1\ne0\). All ratios \(c_j/c_1\) are fixed by every
automorphism over \(k\), so belong to \(k\). For a number field,
the fixed field is \(k\): an algebraic element outside \(k\) has a
different conjugate, and a transcendental element may be moved in
a transcendence basis, in either case extending the isomorphism
to the algebraic closure. Write \(b'=c_1 b\) with \(b\in A\).
Then \(u=a b^4\) with \(a=c_1^4\in k^\times\), because
\(1\otimes u=c_1^4(1\otimes b^4)\) and \(u,b^4\in A\) are nonzero.

Fix a generic frame on each component \(Z\), and absorb the
frame character into \(b\). This gives a rational element \(b_Z\)
in its function field. In another frame with character \(j\),
\((j b_Z)^4\) is a nonzero coefficient-field multiple of the local
unit representing \(K_4\). Both \(j b_Z\) and its inverse satisfy
monic equations over the smooth chart ring. Normality therefore
makes them regular there. The rational rules of opposite weights
glue to inverse regular rules. Equality of their fourth powers with
\(K_4\) up to \(a_Z\in k^\times\) can be checked at the generic
frame. This proves \(e_Z\rho^4=a_Ze_ZK_4\) once the continuous
root and the chartwise unit above are given.
\footnote{\href{https://github.com/ChristianKudeba/alpoge-rh-verification/blob/alpoge-source-7cbf7a0c/src/Development/AbelianVarietyTheory/Proofs/IwasawaDecomposition/ContinuousOnKeySectionComplexValue.lean}{Continuity from the adapted period chart and the theta scalar};
\href{https://github.com/ChristianKudeba/alpoge-rh-verification/blob/alpoge-source-7cbf7a0c/src/Development/AbelianVarietyTheory/Proofs/IwasawaDecomposition/ExistsKeySectionUnif.lean}{descent of the power up to a base-field scalar};
\href{https://github.com/ChristianKudeba/alpoge-rh-verification/blob/alpoge-source-7cbf7a0c/src/Development/AbelianVarietyTheory/Proofs/DivisorProofs/ExistsKeySectionFourthRootPerComponent.lean}{the componentwise rational root and normality argument}.}

\subsection{The integral scalar in two different Lie frames}

Let \(T\) be a complete discrete valuation ring of characteristic-zero
fraction field \(F\). Fix an integral framed point and a prime
\(w\) of \(k\), with the frame valuation restricting to \(e v_w\),
\(e>0\). Take the split eigendirections \(u_{\tau i}\) furnished by
the point, \(1\le i\le m_\tau\). For each \(\tau\) they span the
entire integral \(\tau\)-eigenspace in the tangent module. They
form a basis collectively on the fraction field. Write \(U\) for
the square matrix whose rows are these vectors in an integral Lie
basis and set \(D=\det U\). The good-model theorem for the theta
scalar applies to an integral generating theta section in an
integral Lie basis. The assumptions in the source application
include \(p_0+q_0\ge1\) and the level predicate
\srcid{GoodAt} at every rational prime. The primitive scalar is
then a unit \(\eta\in T^\times\); its comparison with the eigenframe
gives
\[
                            K_4(y)=D^4\eta .
\]
In particular, in this equation unitness refers to \(\eta\).
The eigenframe can have positive colength at ramified primes.

For clarity on the scope of the good-model input, the release
uses the invertibility of the theta pushforward, its model
uniformization, the arc property, and the split-fiber theorem. On
the normal affine chart it starts with a unit on the dense
generic open. A rational function on a normal noetherian affine
domain extends if it belongs to all its height-one valuation rings.
The proof applies this test to both the function and its reciprocal.
At the boundary the arc comparison supplies bounded complex values
for both; at the good-reduction tests the good-reduction unit theorem
does the same algebraically. The local extensions multiply to one
by density. The resulting chart unit pulls back to the trait, where
a change between two integral generating frames multiplies it by an
eighth power of a unit and the inverse fourth power of a determinant
unit. The arc comparison and good-reduction assertions are separate
geometric inputs. An arbitrary chart centered at a bad level prime
does not acquire these hypotheses from the field-valued connection.

\subsection{The different controls the colength}

Put \(\mathfrak D=\mathfrak D_{E/\mathbb Q}\).
For each \(\tau\) choose a nonzero
\(\delta_\tau\in\mathfrak D\) such that
\[
    v_w(\tau(\delta_\tau))
      =\operatorname{ord}_w(\tau(\mathfrak D)\mathcal O_k).
\]
Such an element exists: the extended ideal is contained in the
power of \(w\) on the right and fails containment in the next power.
If all images of elements of \(\mathfrak D\) lay in the next power,
so would the generated ideal.

Let \(\omega_j\) be a \(\mathbb Z\)-basis of \(\mathcal O_E\) and
\(\omega_j^*\) its trace dual. By definition of the inverse different
\(\delta_\tau\omega_j^*\in\mathcal O_E\). Let \(A(a)\) be the matrix
for \(a\) in the integral Lie basis. Set
\[
   d_\tau=\tau(\delta_\tau),\qquad
   P_\tau=\sum_j \tau(\delta_\tau\omega_j^*)A(\omega_j).
\]
All the coefficients on the right are integral in \(T\), via the
given embedding. Over \(F\) the \(d_\tau^{-1}P_\tau\) are the
character projectors. Indeed
\(\sum_j\tau(\omega_j^*)\sigma(\omega_j)=\delta_{\tau\sigma}\).
To verify this, form the square embedding matrices of the two
bases. Their pairing under the trace is the identity, and inverting
in the opposite order gives exactly this relation. The embeddings
indexed here are all embeddings of \(E\), whose images belong to
\(F\). The relation acts on each eigenspace as claimed and yields
\(\sum_\tau d_\tau^{-1}P_\tau=I\).

Every column of \(P_\tau\) is an integral eigenvector of character
\(\tau\). Resolve it in the generating vectors \(u_{\tau i}\).
Stacking the integral coefficients gives a matrix \(G\) whose
\((j,(\tau,i))\)-entry is the coefficient of \(u_{\tau i}\) in
column \(j\) of \(P_\tau\). The transpose of the projector sum is
then
\[
  G\,\operatorname{diag}(d_\tau^{-1}
           \text{ repeated }m_\tau\text{ times})\,U=I .
\]
Taking determinants yields the integral identity
\(\det G\cdot D=\prod_\tau d_\tau^{m_\tau}\).
Thus
\[
  0\le \frac{v_T(K_4(y))}{e}
   =\frac{4v_T(D)}{e}
   \le 4\sum_\tau m_\tau
          \operatorname{ord}_w(\tau(\mathfrak D)\mathcal O_k).
\]
The estimate is zero wherever all the displayed different ideals
are units. If the fourth root on \(Z\) uses \(a_Z\), then the
weight-one section \(e_Z\rho^{-1}\), at any nonzero evaluation
of \(e_Z\), has normalized order in
\[
 \left[-\frac{v_w(a_Z)}4
       -\sum_\tau m_\tau\operatorname{ord}_w(\tau(\mathfrak D)\mathcal O_k),
       \ -\frac{v_w(a_Z)}4\right].
\]
This is the vertical contribution; on a horizontal trait the
coefficient field lies in \(T\) and the valuation comparison in
Section~\ref{sec:jet-connection} directly gives order zero.
\footnote{\href{https://github.com/ChristianKudeba/alpoge-rh-verification/blob/alpoge-source-7cbf7a0c/src/Development/AbelianVarietyTheory/Proofs/DivisorProofs/ColengthDvdOfNeZero.lean}{Trace-dual projectors and the determinant divisibility};
\href{https://github.com/ChristianKudeba/alpoge-rh-verification/blob/alpoge-source-7cbf7a0c/src/Development/AbelianVarietyTheory/Proofs/DivisorProofs/ThetaKeySectionValuationBoundsDifferent.lean}{the exact normalized vertical bounds};
\href{https://github.com/ChristianKudeba/alpoge-rh-verification/blob/alpoge-source-7cbf7a0c/src/Development/AbelianVarietyTheory/Proofs/DivisorProofs/ThetaKeyScalarIsUnitOnIntegralModels.lean}{the model hypotheses and height-one extension}.}
